\newpage
\section{问题重述}

\subsection{问题背景}
城市交通信号控制是缓解交通拥堵的有效手段。传统固定配时方案无法适应动态变化的交通流，需要基于实时交通状态的自适应控制策略。多路口协同控制可以显著减少整体延误，提升路网通行能力。

\subsection{问题提出}
\textbf{问题一}：分析路口交通流数据，提取交通流的时空演化特征，识别高峰期和拥堵传播模式。

\textbf{问题二}：建立多路口协同信号控制优化模型，以总延误时间最小化为目标，考虑绿波带协调和相位差约束。

\textbf{问题三}：设计基于强化学习的自适应信号控制Agent，在动态交通环境下实现实时配时优化。

\subsection{研究意义}
为城市智能交通信号控制系统设计提供算法支撑。
